\begin{mistake}{2026-04-14}{04}

\begin{problemcontent}
已知 \(A,B\) 均为 \(n\) 阶非零矩阵，且秩 \(r(A) = r(B)\)，则必有（\quad）。

(A) \(r(A, B) = r(A)\) 

(B) \(r(A, B) = 2r(A)\)

(C) \(r(A, B) \leqslant 2r(A)\) 

(D) \(r(A - B) = 0\)
\end{problemcontent}

\begin{solutioncontent}
正确选项为 (C)。

根据分块矩阵的秩的性质，拼接矩阵 \((A, B)\) 的列向量组由 \(A\) 和 \(B\) 的列向量组组合而成，其极大无关组的大小不会超过两者极大无关组大小之和，即：
\[ r(A, B) \leqslant r(A) + r(B) \]
将已知条件 \(r(A) = r(B)\) 代入上式，即可得到：
\[ r(A, B) \leqslant r(A) + r(A) = 2r(A) \]
因此 (C) 选项必然成立。

反例排除法：
设 \(A = \begin{pmatrix} 1 & 0 \\ 0 & 0 \end{pmatrix}, B = \begin{pmatrix} 0 & 0 \\ 0 & 1 \end{pmatrix}\)，满足 \(r(A)=r(B)=1\)。
此时 \((A, B) = \begin{pmatrix} 1 & 0 & 0 & 0 \\ 0 & 0 & 0 & 1 \end{pmatrix}\)，则 \(r(A, B) = 2 \neq r(A)\)，排除 (A)。
此时 \(A - B = \begin{pmatrix} 1 & 0 \\ 0 & -1 \end{pmatrix}\)，\(r(A-B) = 2 \neq 0\)，排除 (D)。

另设 \(A = \begin{pmatrix} 1 & 0 \\ 0 & 0 \end{pmatrix}, B = \begin{pmatrix} 1 & 0 \\ 0 & 0 \end{pmatrix}\)，满足 \(r(A)=r(B)=1\)。
此时 \((A, B) = \begin{pmatrix} 1 & 0 & 1 & 0 \\ 0 & 0 & 0 & 0 \end{pmatrix}\)，则 \(r(A, B) = 1 \neq 2r(A)\)，排除 (B)。
\end{solutioncontent}

\mistakeknowledge{
\begin{itemize}
 \item \textbf{核心定理}：分块矩阵列拼接的秩不等式 \(r(A, B) \leqslant r(A) + r(B)\)。
 \item \textbf{易错点}：混淆“秩相等”与“列空间相同”。\(r(A)=r(B)\) 推不出 \(B\) 的列向量能由 \(A\) 线性表出，故不能轻率选 (A)。
 \item \textbf{关联知识}：\(r(A, B) = r(A)\) 成立的充要条件是矩阵方程 \(AX=B\) 有解。
\end{itemize}}

\end{mistake}
